\chapter{Exact weak-C2 source complement after the small-gap admission}
\label{component:pe2fc3cbe3819}
\noindent\textbf{Component }\path{pe2fc3cbe3819}\par
{\small\noindent Source: \path{R13/advances/revision13_weak_local/frontier_audit.tex}\par}
{\small\noindent SHA256: \path{e2fc3cbe38194b492857ca85f59e74674af370627b314af72c93af4b4e1b8c64}\par}
{\small\noindent Inventory use: R9 source/remote extension WF4-WF7 and exact ordered source table\par}
\medskip
\section{Exact weak-C2 source complement after the small-gap admission}
\label{pe2fc3cbe3819:sec:r13-weak-frontier}

This audit uses the original source alternatives (WII1)--(WII6)
and precisely the proved WA, WE, SL, and WC admissions.
The finalized WC proof and its checker are unchanged.

\paragraph{Common quantified source domain.}
Work in the canonical reflected strong-$C$/weak-$\alpha$ C2 II
orientation. Retain the original physical source and stationary
contact, every strict source gap, $n\ge0$, $d=1$,
$B,z,y>0$ infinitesimal, and the original positive-infinitesimal
total-scale intersection. Every row below is further intersected
with the complement of every inherited admission and every exact
symmetry image. No row is asserted nonempty.

Use consistent notation
\[
 a=z/B,\quad t=y/B,\quad h=H/B,\quad l=1-h=L/B,\quad
 e=1-r,\quad g=G/B.
\]
The positive bounded quantities $h,l,r,e,g,j$ each have type
$\mathrm I$ (infinitesimal) or $\mathrm U$ (positive bounded unit).
The quantities $a,t>0$ can additionally have type
$\mathrm O$ (unbounded). Retain exactly
\[
 g=1-ae+hz>0,\qquad j=1+zr-l/t>0,\qquad B-zyr>0
\]
and every other original physical/contact condition. Shape entries
are not independent choices.

Some useful exact implications are
\[
 \begin{gathered}
 ae<1+hz,\qquad e\in\mathrm U\Longrightarrow a\text{ bounded},
 \qquad a\in\mathrm O\Longrightarrow e\in\mathrm I,\\
 a\in\mathrm I\Longrightarrow g\sim1,\qquad
 g\in\mathrm I\Longrightarrow ae\sim1.
 \end{gathered}                                                \tag{WF1}
\]
Thus infinitesimal $a$ is impossible when $g$ is infinitesimal.
For $t$ a unit and $h$ infinitesimal, $l$ is a unit; $j$ a unit
implies $t_0>1$, whereas $j$ infinitesimal implies $t_0=1$.
The directed strong-$C$ condition also retains the exclusion from
(WII6): $r,g,H/y,j$ cannot all be infinitesimal.

\paragraph{Exactly what the current admissions remove.}
Call $t,g,j$ units and $h$ infinitesimal the regular-gap $A$ opening.
On this source $l$ is a unit. If $e$ is a unit, $a$ is bounded by
(WF1). A unit $r$ is covered by WA for unit $a$ and by WE for
infinitesimal $a$; infinitesimal $r$ is covered by SL. Therefore
every $e$-unit source on this opening is admitted, with no omitted
$r$-type. If $e$ is infinitesimal, $r$ is a unit and WC removes
exactly the unit-$a$ type. The exact remainder on this opening is
\[
 \boxed{\quad e\in\mathrm I,\qquad
             a\in\mathrm I\ \text{or}\ a\in\mathrm O.\quad}      \tag{WF2}
\]

\paragraph{An ordered disjoint residual partition.}
With the common domain and all standing exclusions imposed, the
remaining source shapes are the following mutually disjoint rows:
\[
\begin{array}{c|l}
 R1&t\in\mathrm I\\
 R2&t\in\mathrm O\\
 R3&t,h\in\mathrm U,\quad l\in\mathrm I\\
 R4&t,h,l\in\mathrm U,\quad j\in\mathrm I\\
 R5&t,h,l,j\in\mathrm U,\quad r\in\mathrm I\\
 R6&t,h,l,j,r\in\mathrm U,\quad g\in\mathrm I\\
 R7&t\in\mathrm U,\ h,g\in\mathrm I,\quad
                     \text{excluding }r,j\in\mathrm I\\
 R8&t,g\in\mathrm U,\quad h,j\in\mathrm I\\
 R9&t,g,j\in\mathrm U,\quad h,e,a\in\mathrm I\\
 R10&t,g,j\in\mathrm U,\quad h,e\in\mathrm I,\ a\in\mathrm O .
\end{array}                                                    \tag{WF3}
\]
In R7 the excluded $r,j$-intersection is exactly the already
removed two-weak case, because $H/y=h/t$ is infinitesimal.

No $a$-type is silently discarded. R1--R4 and R8 retain all three
positive $a$-types subject to (WF1) and the original source
constraints. R5 has bounded $a$, hence type $\mathrm I$ or
$\mathrm U$. R6--R7 have unit or unbounded $a$, with $ae\sim1$.
R9 and R10 specify the remaining $a$-type explicitly.

The unbounded original coefficients supplied by (WII3) are:
$C,D$ in R1; $C,A$ in R2; $C$ in R3; $D$ in R4;
$n$ in R5; $m$ in R6; $A,m$ in R7; $A,D$ in R8;
and $A$ in R9--R10. Additional $n,m,D$ openings remain governed
by the exact $r,g,j$ conditions. In R1, $A$ is weak; in R2, $D$
is weak. Two unbounded adjacent coefficients alone are not an
invocation of a different dominant-event theorem.

For exhaustiveness, first split $t$ into its three types. At unit
$t$, split $h$ into infinitesimal or unit. For unit $h$, split $l$,
then $j$, then $r$, then $g$. The all-unit remainder has
positive-unit total scale by (WII4), hence is absent from the
original residual intersection. These are R3--R6.
For infinitesimal $h$, split $g$ and then $j$, giving R7--R8.
With both gaps units, (WF2) leaves exactly R9--R10. This proves
the partition, without asserting nonemptiness after all inherited
admissions and symmetry images are removed.

\subsection{A remote reduction for the simultaneous small-$e$, small-$a$ opening}
\label{pe2fc3cbe3819:sec:r13-joint-a-e-remote}

\begin{lemma}
On R9, every nonlocal independent testing chart is finite tensor
SOS, with every original state component and literal projective
endpoint retained. The only remaining testing problem is the full
local parent with both $u,v$ infinitesimal or zero. This is a remote
reduction, not a new all-tests source admission.
\end{lemma}

The exact normalized source (WE1) remains valid without a unit-$a$
assumption. On R9, its polynomials satisfy
\[
 \Delta\sim1,\qquad
 K=(1+a)(h+t)(1+aB)\sim t,\qquad Q\sim1,               \tag{WF4}
\]
where $Q$ is the exact positive polynomial in (BS2). These are
joint reductions: $B,a,h,e$ are infinitesimal, $r=1-e$, and $t$
is a positive unit. The same exact source identity gives
\[
 2r\Delta n=eQ-BK,\qquad
 0<B/e\le Q/K,\qquad 0\le n/e\le Q/(2r\Delta).        \tag{WF5}
\]
In particular $n$ is infinitesimal, including literal $n=0$.

Since $g=1+a(hB-e)\sim1$ and
$j=1+aBr-(1-h)/t\sim1-1/t$ is a positive unit, $t_0>1$.
Writing $c_a=1+aBr$, $S_c=1+B(1-h+ar)$, and
\[
 M_*=(1+aB)(1+ar)(h+t)-a(1+tB)e(1-h),
\]
the remaining exact source formulas yield
\[
\begin{aligned}
 \frac{m}{a^2e}
 &=\frac{(1+tB)c_aS_c}{2g\Delta}\longrightarrow\frac12,\\
 Ah&=\frac{tc_a(1+ar)S_c}{2\Delta}\longrightarrow t_0/2,\\
 C&=\frac{c_aM_*}{2(1-h)\Delta}\longrightarrow t_0/2,\\
 D&=\frac{(1+aB)(1+ar)(1-atBr)S_c}{2tj\Delta}
                                  \longrightarrow\frac1{2(t_0-1)},\\
 \frac{\alpha}{B^2}
 &=\frac{(1+ar)M_*}{2(1-hB)\Delta}\longrightarrow t_0/2.
\end{aligned}                                                     \tag{WF6}
\]
All denominators have positive-unit residues. Hence
$m=O(a^2e)$, $\alpha=O(B^2)$, $A^{-1/2}=\sqrt{h/(Ah)}$ is
infinitesimal, and the complete coefficient conditions actually
used by (GR2)--(GR7) hold:
\[
 n,m,\alpha,A^{-1/2}\in\operatorname{Inf},\qquad
 C_0=t_0/2,\qquad D_0=1/[2(t_0-1)].                  \tag{WF7}
\]
Those exact complete-parent identities contain no inverse $a$.
Their small $m$ packet is retained inside $\sqrt m$, and the
$n$ and $\alpha$ roots permit zero weights without division.
Every printed strict residue determinant is consequently unchanged.
The existing coefficient-bound and homogeneous-endpoint arguments
therefore apply verbatim to every nonlocal test. This proves the
lemma without replaying any unchanged algebraic checker. It does
not import the local WC circuit.

\paragraph{The remaining local obstruction is a normalization issue.}
WC does not yet give the R9 local theorem. Its present scalar
coordinates contain
\[
 s=\sqrt{t/a}\ \text{unbounded},\qquad
 a\partial_U f_C=\alpha/(tB^2)\longrightarrow\tfrac12.
                                                               \tag{WF8}
\]
The unit-$s$ hypothesis and bounded finite coefficient arrays of
the three-circuit lift therefore fail at the new source boundary.
There is no license to call $U/a$ small: the source and testing
fields are independent. In the WC reciprocal frame, the bound on
$n/m$ is also lost; (WF5)--(WF6) give only $n/m=O(a^{-2})$,
so they do not control $\sqrt{n/m}(u/v)^2$ uniformly.

The older WE-rescaled WA frame has the same separate failure:
scaling its angular root by $\sqrt e$ leaves nonzero-residue
$J_s/(a\sqrt e)$ and $f_s/(a\sqrt e)$ coefficients. These are
limitations of the stated frames, not non-SOS obstructions.

The smallest decisive new lemma is a contact lift with the
reciprocal source slope $c=\sqrt{a/t}$ infinitesimal and source
Bob radius $\sqrt y$ infinitesimal, retaining the whole centered
mass $q_s=nr^2+(m+1/2)e^2$. It must give explicit finite conductors
for the $U/a$ opening and retain $\sqrt n\,uZ$, or an exact
contact-centered replacement, as an additional weighted root
instead of dividing it by $\sqrt m$. It must include both full
soft states, every normal, $n=0$, arbitrary $h/a$, and all
independent testing ratios. No such lemma is proved here.

R10 is separate: $a$ is unbounded, $ae$ is bounded, and $g$ remains
a positive unit. Normalizing by the source product $ae$ is a
natural next calculation, but neither the unit-$a$ WC theorem nor
the R9 remote reduction supplies that source-domain admission.
